\documentclass[12pt]{report}

% ================================================================
% FRAMEWORK ROOT
% Prefix path menuju akar framework. Kosong bila dokumen berada
% di root; untuk contoh (examples/<jenis>/main.tex) gunakan "../../".
% Seluruh \input internal framework memakai \frameworkRoot.
% ================================================================

\newcommand{\frameworkRoot}{}

% ================================================================
% IDENTITAS DOKUMEN
% Ganti semua nilai di bawah ini dengan identitas penelitian Anda.
% Nilai yang ada saat ini hanya contoh (placeholder).
% ================================================================

\newcommand{\docTitle}{
    Judul Penelitian Anda
}

\newcommand{\docAuthor}{
    Nama Lengkap Penulis
}

\newcommand{\Pembimbing}{

    \begin{tabular}{c}
        Nama Pembimbing 1 \\
        Nama Pembimbing 2 \\
        Nama Pembimbing 3 \\
    \end{tabular}
}

% ================================================================
% INSTITUSI
% ================================================================

\newcommand{\Program}{
    PROGRAM STUDI
}

\newcommand{\Department}{
    DEPARTEMEN
}

\newcommand{\Faculty}{
    FAKULTAS
}

\newcommand{\Hospital}{
    RUMAH SAKIT
}

\newcommand{\University}{
    UNIVERSITAS
}

\newcommand{\Location}{
    KOTA
}

\newcommand{\SubmissionYear}{
    2026
}

% ================================================================
% METADATA PDF
% ================================================================

\newcommand{\docSubject}{
    Proposal/Laporan Penelitian
}

\newcommand{\docKeywords}{
    Kata kunci 1, Kata kunci 2, Kata kunci 3
}

\usepackage[a4paper,left=4cm,right=3cm,top=3cm,bottom=3cm]{geometry}

\usepackage{fontspec}
\setmainfont{Times New Roman}

\usepackage[indonesian]{babel}
\usepackage{setspace}
\usepackage{graphicx}
\usepackage{booktabs}
\usepackage{array}
\usepackage{titlesec}
\usepackage{caption}
\usepackage{csquotes}
\usepackage{microtype}
\usepackage{amsmath}
\usepackage{chngcntr}

\usepackage[
backend=biber,
style=vancouver,
sorting=none,
citestyle=numeric-comp,
language=english
]{biblatex}

\DeclareLanguageMapping{indonesian}{english}
\addbibresource{bibliography/references.bib}
% ================================================================
% SETTINGS
% Induk modul pengaturan tampilan. Setiap modul memiliki satu
% tanggung jawab. Tambahkan modul baru dengan meng-input-nya di sini.
%
% Path modul diawali \frameworkRoot — makro yang didefinisikan oleh
% entry point (main.tex: kosong; contoh: ../../). Dengan ini framework
% dapat dikompilasi dari direktori mana pun.
% ================================================================

% ================================================================
% SETTINGS — SPACING
% Jarak baris dan paragraf.
% ================================================================

\onehalfspacing

\setlength{\parindent}{1.27cm}
\setlength{\parskip}{0pt}

% ================================================================
% SETTINGS — TYPOGRAPHY
% Nama-nama bagian dokumen (daftar isi, gambar, tabel, pustaka).
% ================================================================

\renewcommand{\contentsname}{DAFTAR ISI}
\renewcommand{\listfigurename}{DAFTAR GAMBAR}
\renewcommand{\listtablename}{DAFTAR TABEL}
\renewcommand{\bibname}{DAFTAR PUSTAKA}

\renewcommand{\figurename}{Gambar}
\renewcommand{\tablename}{Tabel}

% ================================================================
% SETTINGS — CHAPTER
% Format judul BAB dan penomorannya (Romawi).
% ================================================================

\renewcommand{\thechapter}{\Roman{chapter}}

\titleformat{\chapter}[display]
{\centering\bfseries\fontsize{16pt}{20pt}\selectfont}
{BAB~\thechapter}
{0.5em}
{\MakeUppercase}

\titlespacing*{\chapter}
{0pt}
{0pt}
{20pt}

% ================================================================
% SETTINGS — SECTION
% Format judul section, subsection, dan subsubsection.
% ================================================================

\renewcommand{\thesection}
{\arabic{chapter}.\arabic{section}}

\titleformat{\section}
{\bfseries\fontsize{12pt}{18pt}\selectfont}
{\thesection}
{1em}
{}

\titlespacing*{\section}
{0pt}
{18pt}
{6pt}

\renewcommand{\thesubsection}
{\thesection.\arabic{subsection}}

\titleformat{\subsection}
{\bfseries\fontsize{12pt}{18pt}\selectfont}
{\thesubsection}
{1em}
{}

\titlespacing*{\subsection}
{0pt}
{18pt}
{6pt}

\renewcommand{\thesubsubsection}
{\thesubsection.\arabic{subsubsection}}

\titleformat{\subsubsection}
{\bfseries\fontsize{12pt}{18pt}\selectfont}
{\thesubsubsection}
{1em}
{}

\titlespacing*{\subsubsection}
{0pt}
{18pt}
{6pt}

% ================================================================
% SETTINGS — NUMBERING
% Penomoran gambar dan tabel per bab.
% ================================================================

\numberwithin{figure}{chapter}
\numberwithin{table}{chapter}

\renewcommand{\thefigure}
{\arabic{chapter}.\arabic{figure}}

\renewcommand{\thetable}
{\arabic{chapter}.\arabic{table}}

% ================================================================
% SETTINGS — CAPTION
% Format keterangan gambar dan tabel.
% ================================================================

\captionsetup{font=small,
  labelfont=bf,
  textfont=normalfont,
  labelsep=space,
  justification=justified,
  singlelinecheck=false,
  skip=8pt}

% ================================================================
% SETTINGS — TOC / LOF / LOT
% Format daftar isi, daftar gambar, dan daftar tabel.
% ================================================================

\makeatletter

\renewcommand{\tableofcontents}{
    \chapter*{
        \centering
        \bfseries
        \fontsize{16pt}{20pt}\selectfont
        DAFTAR ISI
    }

    \@starttoc{toc}
}

\renewcommand{\listoffigures}{
    \chapter*{
        \centering
        \bfseries
        \fontsize{16pt}{20pt}\selectfont
        DAFTAR GAMBAR
    }

    \@starttoc{lof}
}

\renewcommand{\listoftables}{
    \chapter*{
        \centering
        \bfseries
        \fontsize{16pt}{20pt}\selectfont
        DAFTAR TABEL
    }

    \@starttoc{lot}
}

\makeatother

% ================================================================
% SETTINGS — BIBLIOGRAPHY
% Format heading daftar pustaka.
% ================================================================

\defbibheading{bibliography}{
\chapter*{
\centering
\bfseries
\fontsize{16pt}{20pt}\selectfont
DAFTAR PUSTAKA
}
\addcontentsline{toc}{chapter}{DAFTAR PUSTAKA}
}

% ================================================================
% SETTINGS — HYPERLINK
% Konfigurasi hyperref (tautan internal PDF) dan metadata PDF.
% Metadata dokumen diambil dari preamble/metadata.tex (\doc*).
% ================================================================

\usepackage[
colorlinks=true,
linkcolor=black,
citecolor=black,
urlcolor=black,
unicode=true,
pdftitle={\docTitle},
pdfauthor={\docAuthor},
pdfsubject={\docSubject},
pdfkeywords={\docKeywords}
]{hyperref}


% ================================================================
% COMMANDS
% Perintah (command) khusus yang dapat dipakai ulang.
% Semua nilai dokumen berasal dari preamble/metadata.tex (\doc*).
% ================================================================

\newcommand{\HRule}{\rule{\linewidth}{0.5mm}}

% ================================================================
%% FIGURE
%% Sisipkan gambar dengan caption dan label dalam satu perintah.
%%
%% Pemakaian:
%%   \Figure{figures/logo_institusi.png}{Keterangan gambar.}{fig:logo}
%%
%% Argumen: [1] file gambar, [2] caption, [3] label.
%% ================================================================

\NewDocumentCommand{\Figure}{m m m}{
  \begin{figure}[ht]
  \centering
  \includegraphics[width=0.8\linewidth]{#1}
  \caption{#2}
  \label{#3}
  \end{figure}
}

% ================================================================
%% TABLE
%% Sisipkan tabel dengan caption dan label dalam satu perintah.
%% Isi tabel ditulis dalam argumen kedua (gaya booktabs).
%%
%% Pemakaian:
%%   \Table{tab:contoh}{%
%%     \begin{tabular}{ll}
%%       \toprule
%%       A & B \\
%%       \bottomrule
%%     \end{tabular}
%%   }{Keterangan tabel.}
%%
%% Argumen: [1] label, [2] isi tabel, [3] caption.
%% ================================================================

\NewDocumentCommand{\Table}{m m m}{
  \begin{table}[ht]
  \centering
  #2
  \caption{#3}
  \label{#1}
  \end{table}
}

% ================================================================
%% EQUATION
%% Sisipkan persamaan bernomor dengan label.
%%
%% Pemakaian:
%%   \Equation{eq:contoh}{E = mc^2}
%%
%% Argumen: [1] label, [2] isi persamaan.
%% ================================================================

\NewDocumentCommand{\Equation}{m m}{
  \begin{equation}
  #2
  \label{#1}
  \end{equation}
}

% ================================================================
%% APPENDIX
%% Mulai bagian lampiran. Membaca label optional untuk penamaan.
%%
%% Pemakaian:
%%   \Appendix{Lampiran}        % bagian lampiran, judul "Lampiran"
%%   \Appendix*                 % lampiran tanpa nomor bab baru
%% ================================================================

\NewDocumentCommand{\Appendix}{s m}{
  \IfBooleanTF{#1}
    {\chapter*{#2}}
    {\chapter{#2}}
}

% ================================================================
%% SOURCE
%% Catatan sumber untuk gambar/tabel (ditempatkan setelah lingkungan).
%%
%% Pemakaian:
%%   \Figure{figures/x.png}{Keterangan.}{fig:x}
%%   \Source{Diadaptasi dari: Author (2020).}
%% ================================================================

\NewDocumentCommand{\Source}{m}{
  \par\smallskip
  \begin{flushleft}
  \small #1
  \end{flushleft}
}

% ================================================================
%% MAKECOVER
%% Halaman sampul (titlepage) yang dapat dipakai ulang.
%% Membaca seluruh identitas dari metadata.tex.
%%
%% Pemakaian:
%%   \makecover                          % judul dari \docTitle
%%   \makecover[Judul Pengganti]         % override judul
%%   \makecover[][Skripsi]               % jenis dokumen (proposal/skripsi/tesis/disertasi/laporan)
%%
%% Contoh jenis dokumen: Proposal, Skripsi, Tesis, Disertasi,
%% Laporan Akhir, dsb. Biarkan argumen ke-2 kosong untuk mengabaikan.
%% ================================================================

\makeatletter

\NewDocumentCommand{\makecover}{ o o }{%
  \begin{titlepage}

  \thispagestyle{empty}

  \begin{center}

  % JUDUL
  {\bfseries
  \fontsize{16pt}{20pt}\selectfont
  \MakeUppercase{\IfValueTF{#1}{#1}{\docTitle}}
  }

  % OPTIONAL: SUBJUDUL
  \ifx\docSubtitle\@empty\relax\else
  {\fontsize{14pt}{18pt}\selectfont
  \docSubtitle
  }
  \fi

  % OPTIONAL: JENIS DOKUMEN
  \IfValueT{#2}{%
  \vspace{1cm}
  {\fontsize{12pt}{18pt}\selectfont
  #2
  }
  }

  \vspace{1.5cm}

  % LOGO (path dari \docLogo)
  \includegraphics[width=4cm,height=4cm,keepaspectratio]{\docLogo}

  \vspace{2cm}

  % PENULIS
  {\fontsize{12pt}{18pt}\selectfont
  Disusun oleh
  }

  \vspace{0.5cm}

  {\bfseries
  \fontsize{12pt}{18pt}\selectfont
  \docAuthor
  }

  \vspace{1.5cm}

  % PEMBIMBING
  {\fontsize{12pt}{18pt}\selectfont
  Telah diperiksa dan disetujui oleh : \\
  }

  \vspace{0.5cm}

  {\bfseries
  \fontsize{12pt}{18pt}\selectfont
  \docSupervisor
  }

  \vfill

  % INSTITUSI
  {\bfseries
  \fontsize{12pt}{18pt}\selectfont

  \docProgram

  %\docDepartment

  %\docFaculty

  \docHospital

  \docLocation

  \docYear

  }

  \end{center}

  \end{titlepage}

  \clearpage
}

\makeatother


% ================================================================
% THEME
% Pilih tema tampilan: classic (default), modern, minimal, book.
% ================================================================

\usetheme{classic}

% ================================================================
% INCLUDEONLY
% Kompilasi cepat: hanya memuat bab tertentu, misalnya
%   \includeonly{chapters/01_pendahuluan, chapters/02_tinjauan_pustaka}
% Catatan: nomor halaman/sitasi bisa berubah karena file .aux lain
% tidak ikut diproses. Gunakan untuk editing cepat, lalu build penuh.
% ================================================================

%\includeonly{chapters/01_pendahuluan}

\begin{document}

\begin{titlepage}

\thispagestyle{empty}

\begin{center}

%%%%%%%%%%%%%%%%%%%%%%%%%%%%%%%%%%%%%%%%%%%%%%%%%%%%%%%%%%%%
%% JUDUL
%%%%%%%%%%%%%%%%%%%%%%%%%%%%%%%%%%%%%%%%%%%%%%%%%%%%%%%%%%%%

{\bfseries
\fontsize{16pt}{20pt}\selectfont
\MakeUppercase{\docTitle}
}

\vspace{2cm}

%%%%%%%%%%%%%%%%%%%%%%%%%%%%%%%%%%%%%%%%%%%%%%%%%%%%%%%%%%%%
% LOGO
% Path logo diambil dari \docLogo (preamble/metadata.tex).
% Simpan berkas logo di folder figures/.
%%%%%%%%%%%%%%%%%%%%%%%%%%%%%%%%%%%%%%%%%%%%%%%%%%%%%%%%%%%%

\includegraphics[width=4cm,height=4cm,keepaspectratio]{\docLogo}

\vspace{2cm}

%%%%%%%%%%%%%%%%%%%%%%%%%%%%%%%%%%%%%%%%%%%%%%%%%%%%%%%%%%%%
%% PENULIS
%%%%%%%%%%%%%%%%%%%%%%%%%%%%%%%%%%%%%%%%%%%%%%%%%%%%%%%%%%%%

{\fontsize{12pt}{18pt}\selectfont
Disusun oleh
}

\vspace{0.5cm}

{\bfseries
\fontsize{12pt}{18pt}\selectfont
\docAuthor
}

\vspace{1.5cm}

%%%%%%%%%%%%%%%%%%%%%%%%%%%%%%%%%%%%%%%%%%%%%%%%%%%%%%%%%%%%
%% PEMBIMBING
%%%%%%%%%%%%%%%%%%%%%%%%%%%%%%%%%%%%%%%%%%%%%%%%%%%%%%%%%%%%

{\fontsize{12pt}{18pt}\selectfont
Telah diperiksa dan disetujui oleh : \\
}

\vspace{0.5cm}

{\bfseries
\fontsize{12pt}{18pt}\selectfont
\docSupervisor
}

\vfill

%%%%%%%%%%%%%%%%%%%%%%%%%%%%%%%%%%%%%%%%%%%%%%%%%%%%%%%%%%%%
%% INSTITUSI
%%%%%%%%%%%%%%%%%%%%%%%%%%%%%%%%%%%%%%%%%%%%%%%%%%%%%%%%%%%%

{\bfseries
\fontsize{12pt}{18pt}\selectfont

\docProgram

%\docDepartment

%\docFaculty

\docHospital

\docLocation

\docYear

}

\end{center}

\end{titlepage}

\clearpage
\chapter*{Abstrak}
Penelitian deskriptif ini bertujuan untuk memberikan gambaran umum mengenai topik yang diteliti. Bagian ini dapat diisi dengan ringkasan masalah, metode, hasil, dan kesimpulan penelitian secara singkat.

\vspace{1em}
\textbf{Kata kunci:} kata kunci 1, kata kunci 2, kata kunci 3

\cleardoublepage

\pagenumbering{roman}

\tableofcontents
\addcontentsline{toc}{chapter}{Daftar Isi}
\cleardoublepage

\listoffigures
\addcontentsline{toc}{chapter}{Daftar Gambar}
\cleardoublepage

\listoftables
\addcontentsline{toc}{chapter}{Daftar Tabel}
\cleardoublepage

\pagenumbering{arabic}


\chapter{Pendahuluan}

% ================================================================
% PANDUAN PENULISAN BAB I — PENDAHULUAN
% Bab ini menjelaskan alasan dan tujuan penelitian Anda. Tuliskan
% secara naratif, padat, dan didukung sitasi. Perhatikan urutan
% bagian berikut.
% ================================================================

\section{Latar Belakang}

% Tuliskan 4--7 paragraf dengan alur "funnel" (dari umum ke khusus):
%   1) Buka dengan masalah besar/topik yang relevan dengan penelitian.
%   2) Persempit ke masalah spesifik yang akan diteliti.
%   3) Tegaskan urgensi masalah: data prevalensi/insidensi dan dampaknya.
%   4) Tunjukkan kesenjangan (research gap) dari penelitian sebelumnya.
%   5) Akhiri dengan pernyataan singkat maksud penelitian (research aim).
% Setiap pernyataan faktual harus disertai sitasi, contoh:
% \supercite{contohReferensi1,contohReferensi2}
%
% Contoh pola kalimat:
%   "X merupakan masalah yang ...\supercite{ref}. Prevalensinya dilaporkan
%    mencapai ...\supercite{ref}. Namun, data populasi lokal masih terbatas.
%    Oleh karena itu, penelitian ini bertujuan untuk ..."

\section{Rumusan Masalah}

% Uraikan masalah penelitian dalam bentuk pertanyaan penelitian.
% Rumusan masalah umum ditulis dalam satu kalimat tanya; rumusan
% masalah khusus diturunkan menjadi beberapa butir.
%
% Contoh:
% \begin{enumerate}
%     \item Bagaimana karakteristik demografis (usia dan jenis kelamin) subjek penelitian?
%     \item Bagaimana gambaran klinis dan tata laksana subjek penelitian?
% \end{enumerate}

\section{Tujuan Penelitian}

\subsection{Tujuan Umum}

% Nyatakan tujuan penelitian secara luas, selaras dengan judul penelitian.
% Contoh: "Mengetahui karakteristik klinis pasien ..."

\subsection{Tujuan Khusus}

% Rincikan tujuan umum menjadi butir-butir tujuan yang terukur, dengan
% kata kerja operasional (mengetahui distribusi, mengidentifikasi,
% mengukur, menilai, membandingkan, dll).
% \begin{enumerate}
%     \item Tujuan khusus pertama.
%     \item Tujuan khusus kedua.
% \end{enumerate}

\section{Manfaat Penelitian}

% Jelaskan kegunaan hasil penelitian, misalnya:
% - Manfaat teoretis: pengembangan ilmu/konsep di bidang terkait.
% - Manfaat praktis: masukan bagi klinisi, pelayanan kesehatan,
%   dan perencanaan tata laksana atau kebijakan.
% - Manfaat bagi penelitian selanjutnya: dasar/bahan acuan.

\chapter{Tinjauan Pustaka}

% ================================================================
% PANDUAN PENULISAN BAB II — TINJAUAN PUSTAKA
% Bab ini memuat landasan teoretis yang menunjang penelitian Anda.
% Susun subbab sesuai kebutuhan topik, kemudian kembangkan setiap
% poin menjadi narasi akademis Bahasa Indonesia yang didukung sitasi
% (key referensi sudah tersedia di bibliography/references.bib).
% ================================================================

\section{Definisi dan Konsep Dasar}

% Jelaskan definisi istilah/topik utama penelitian, baik menurut
% literatur maupun klasifikasi/standar yang berlaku. Kemukakan
% mengapa topik tersebut penting untuk dikaji.

\section{Klasifikasi}

% Uraikan sistem klasifikasi yang relevan dengan topik (misal
% berdasarkan etiologi, anatomi, waktu awitan, atau derajat keparahan).
% Tekankan implikasi klinis dari masing-masing klasifikasi terhadap
% tata laksana maupun prognosis.

\section{Epidemiologi}

% Sajikan data epidemiologi topik penelitian:
% - Angka prevalensi/insidensi global dan nasional.
% - Distribusi menurut karakteristik demografis (usia, jenis kelamin).
% - Faktor risiko dan mekanisme yang paling sering ditemukan.
% - Pergeseran beban penyakit antar-region, bila relevan.

\section{Etiologi dan Patofisiologi}

% Jelaskan penyebab dan mekanisme terjadinya penyakit/masalah yang
% diteliti. Bila memiliki beberapa mekanisme, uraikan tiap mekanisme
% dalam subbab/subpoin terpisah agar lebih runtut.

\section{Faktor Risiko dan Prediktor}

% Sebutkan faktor-faktor yang meningkatkan risiko atau menjadi
% prediktor luaran klinis, lengkap dengan sitasi pendukung. Sajikan
% dalam narasi atau tabel perbandingan bila jumlahnya banyak.

\section{Manifestasi Klinis dan Diagnosis}

% - Keluhan dan tanda klinis khas.
% - Pemeriksaan fisik dan penunjang yang digunakan untuk menegakkan
%   diagnosis (mis. pemeriksaan khusus, laboratorium, pencitraan).
% - Urutan/strategi diagnosis dan indikasi pemeriksaan penunjang.

\section{Tata Laksana}

% Uraikan prinsip penatalaksanaan berdasarkan bukti terkini, disusun
% sesuai tahapan (konservatif, medikamentosa, tindakan, bedah).
% Jelaskan indikasi, kelebihan, dan keterbatasan tiap pilihan terapi.

\section{Penelitian Terkait}

% - Bahas studi-studi sebelumnya dari berbagai populasi yang relevan
%   dengan topik penelitian Anda.
% - Tunjukkan persamaan, perbedaan, serta kesenjangan yang masih ada.
% - Tekankan keterbatasan data pada populasi lokal sebagai justifikasi
%   penelitian Anda.

\section{Kerangka Teori, Kerangka Konsep, dan Hipotesis}

% - Kerangka teori: skema alur hubungan antarvariabel berdasarkan teori.
% - Kerangka konsep: skema variabel yang diteliti pada penelitian ini
%   (dapat dibuat sebagai gambar/diagram, lihat contoh di Bab IV).
% - Hipotesis: pernyataan dugaan hasil (diisi bila penelitian bersifat
%   inferensial/analitik; untuk penelitian deskriptif bagian ini
%   dapat disesuaikan).

\section{Ringkasan}

% Rangkum poin-poin kunci tinjauan pustaka sebagai jembatan menuju
% Bab III, misalnya: beban masalah yang tinggi, mekanisme yang
% multifaktorial, serta kesenjangan data pada populasi lokal yang
% melatarbelakangi perlunya penelitian dilakukan.

\chapter{Metode Penelitian}

% ================================================================
% PANDUAN PENULISAN BAB III — METODE PENELITIAN
% Bab ini menjelaskan cara penelitian dilakukan secara rinci sehingga
% dapat direplikasi oleh peneliti lain. Tuliskan secara sistematis
% mengikuti urutan bagian berikut.
% ================================================================

\section{Desain Penelitian}

% Sebutkan desain penelitian Anda, misalnya deskriptif (cross-sectional
% atau longitudinal), case-control, kohort, atau eksperimental. Jelaskan
% singkat alasan pemilihan desain tersebut.

\section{Tempat dan Waktu Penelitian}

% Cantumkan lokasi pengambilan data (misal unit/instalasi rumah sakit,
% klinik, atau laboratorium) serta periode/rentang waktu pengumpulan data.

\section{Populasi dan Sampel Penelitian}

% - Bedakan populasi target dan populasi terjangkau.
% - Jelaskan teknik pengambilan sampel (total sampling, consecutive
%   sampling, purposive sampling, dll).
% - Cantumkan perhitungan besar sampel bila diperlukan.

\subsection{Kriteria Inklusi}

% Daftar karakteristik subjek yang memenuhi syarat untuk diteliti.
% \begin{itemize}
%     \item Kriteria inklusi pertama.
%     \item Kriteria inklusi kedua.
% \end{itemize}

\subsection{Kriteria Eksklusi}

% Daftar karakteristik yang menyingkirkan subjek dari penelitian.
% \begin{itemize}
%     \item Kriteria eksklusi pertama.
%     \item Kriteria eksklusi kedua.
% \end{itemize}

\section{Definisi Operasional}

% Sajikan definisi operasional variabel dalam bentuk tabel: nama
% variabel, definisi, alat ukur, cara ukur, dan hasil ukur/skala.
% Contoh format tabel:
%
% \begin{table}[ht]
% \centering
% \begin{tabular}{p{3.2cm}p{3.6cm}p{2.4cm}p{3cm}p{2.4cm}}
% \toprule
% Variabel & Definisi & Alat Ukur & Cara Ukur & Hasil/Skala \\
% \midrule
% Variabel 1 & ... & ... & ... & ... \\
% Variabel 2 & ... & ... & ... & ... \\
% \bottomrule
% \end{tabular}
% \caption{Definisi operasional variabel}
% \label{tab:definisi}
% \end{table}

\section{Instrumen Penelitian}

% Jelaskan instrumen/alat yang digunakan (kuesioner, formulir pengumpulan
% data, alat pemeriksaan, perangkat lunak, dll) beserta validitas dan
% reliabilitasnya bila relevan.

\section{Prosedur Pengumpulan Data}

% Uraikan langkah pengumpulan data secara berurutan: persiapan,
% seleksi subjek, pemeriksaan/pengisian instrumen, hingga rekapitulasi
% dan verifikasi data.

\section{Analisis Data}

% Sebutkan jenis analisis yang digunakan (deskriptif dan/atau
% inferensial) serta perangkat lunak statistik. Untuk analisis
% inferensial, jelaskan uji statistik yang dipakai sesuai karakteristik
% data (skala, normalitas, distribusi).

\section{Etika Penelitian}

% Jelaskan persetujuan etik (nomor kode etik bila sudah didapat),
% prosedur informed consent, serta jaminan kerahasiaan data subjek.

\chapter{Hasil dan Pembahasan}

% ================================================================
% PANDUAN PENULISAN BAB IV — HASIL DAN PEMBAHASAN
% Sajikan temuan secara objektif pada bagian Hasil, lalu tafsirkan dan
% kaitkan dengan teori serta penelitian lain pada bagian Pembahasan.
% ================================================================

\section{Hasil Penelitian}

% Deskripsikan karakteristik subjek dan hasil analisis sesuai dengan
% tujuan penelitian. Sajikan data berupa narasi yang didukung tabel
% dan/atau gambar. Bagian ini TIDAK boleh berisi penafsiran/interpretasi.

% Contoh gambar (ganti placeholder berikut dengan gambar Anda;
% simpan berkas gambar di folder figures/):
\begin{figure}[ht]
\centering
\fbox{\rule{0pt}{3cm}\rule{0.9\linewidth}{0pt}}
\caption{Tulis keterangan gambar di sini.}
\label{fig:contoh}
\end{figure}

% Alternatif bila gambar sudah tersedia:
% \begin{figure}[ht]
% \centering
% \includegraphics[width=0.8\linewidth]{figures/nama_file_gambar}
% \caption{Tulis keterangan gambar di sini.}
% \label{fig:contoh2}
% \end{figure}

% Contoh tabel:
\begin{table}[ht]
\centering
\begin{tabular}{ll}
\toprule
Variabel & Keterangan \\
\midrule
A & Contoh data A \\
B & Contoh data B \\
\bottomrule
\end{tabular}
\caption{Tulis keterangan tabel di sini.}
\label{tab:contoh}
\end{table}

\section{Pembahasan}

% Interpretasikan hasil penelitian dan kaitkan dengan teori serta hasil
% penelitian sebelumnya. Jelaskan temuan yang mendukung maupun yang
% bertentangan, serta berikan argumen yang logis atas perbedaan tersebut.

\section{Keterbatasan Penelitian}

% Akui kelemahan penelitian (desain, besar sampel, potensi bias) dan
% jelaskan implikasinya terhadap interpretasi serta generalisasi hasil.

\chapter{Penutup}

% ================================================================
% PANDUAN PENULISAN BAB V — PENUTUP
% Bab ini memuat kesimpulan dan saran dari penelitian yang telah
% dilakukan.
% ================================================================

\section{Kesimpulan}

% Jawab setiap rumusan masalah/tujuan penelitian secara ringkas, padat,
% dan jelas, tanpa mengulang seluruh isi hasil. Simpulan harus ditopang
% oleh data yang telah disajikan pada Bab IV.

\section{Saran}

% Berikan rekomendasi berdasarkan hasil penelitian, misalnya bagi:
% - Pelayanan kesehatan: implikasi praktis untuk skrining/tata laksana.
% - Penelitian selanjutnya: perbaikan metodologi atau perluasan cakupan.
% - Institusi: perbaikan sistem pencatatan data, kebijakan, atau dukungan
%   pendanaan penelitian.

\include{\frameworkRoot chapters/06_daftar_pustaka}

\end{document}
